\section{Chiều của phép chuyển, và điều chưa giải quyết}
\label{sec:ch17-chieu-cua-phep-chuyen}

\index{chuỗi suy luận!bốn chỗ đứng của bốn bài|(}%
Bốn mục vừa rồi đọc bốn bài, và mỗi bài đứng ở một chỗ khác nhau so với câu hỏi
của sách. Mục này xếp bốn chỗ ấy lại, và cái xếp ra được là một câu về chiều.

\index{chuỗi suy luận!bốn chỗ đứng khác nhau}%
Bốn chỗ đứng ấy khác nhau ở chỗ mỗi bài làm gì với câu hỏi của
định nghĩa~\ref{def:ch01-do-trung-thuc}. Bài mở đường không trả lời và nói rõ ba
lần rằng nó không trả lời. Khảo sát về mô hình ngôn ngữ cho XAI xếp chuỗi suy
luận vào nhóm không cần hỏi, đặt tiêu chí ấy vào cột định lượng mà không cấp dụng
cụ nào, rồi dẫn một bài đo ngược lại làm chỗ dựa. Khảo sát 2026 thấy câu hỏi, vẽ
hẳn một hình cho nó, gọi tên ba họ dụng cụ, rồi định nghĩa nó thành một câu hỏi
khác. Và bài neo trả lời, cho tám chỉ số, với câu trả lời là các chỉ số ấy chạy
quanh mức ngẫu nhiên.

\index{chuỗi suy luận!chiều thật của phép chuyển}%
Bây giờ tới câu về chiều, và nó là câu của sách. Chỗ đặt chương này trong sách,
cùng với chữ \enquote{chuyển} mà mục~\ref{sec:ch01-luan-de-cua-sach} và
mục~\ref{sec:ch16-ba-cai-gia} đều dùng, gợi rằng lời phê phán của Phần~IV được
mang sang đây. Chiều ấy ngược, và bằng chứng nằm trong chính sách này chứ không
trong bài nào.
Mục~\ref{sec:ch09-doi-tuong-khac-nhau} viết ra tình trạng hiện tại như sau: với
chuỗi suy luận, câu hỏi các chỉ số có đo đúng thứ chúng ghi trên mặt đồng hồ hay
không đã có câu trả lời và câu trả lời là không; với attribution đặc trưng, câu
hỏi ấy vẫn chưa được hỏi bằng thực nghiệm.

Đọc câu ấy như một câu về chiều thì nó nói rằng phép đối chiếu với ground truth
đã được chạy đúng một lần, và nó được chạy ở phía chuỗi suy luận. Bảy chương của
Phần~IV không mang được kết luận nào sang đây, vì kết luận mạnh nhất trong bảy
chương ấy chính là kết luận vừa dẫn, và nó vốn đã ở đây. Cái Phần~IV mang sang
được là một câu hỏi chưa có lời đáp; cái đi ngược lại, từ đây sang Phần~IV, thì
mục~\ref{sec:ch09-doi-tuong-khac-nhau} đã liệt kê thành ba thứ từ tám chương
trước: cách dựng ground truth từ thiết kế nhiệm vụ, hai chẩn đoán phát biểu lại ở
mức cấu trúc, và phép đo mức đồng thuận giữa các chỉ số mà phía attribution chạy
được ngay.

\index{chuỗi suy luận!vì sao chiều lại quan trọng}%
Chiều quan trọng ở ba chỗ. Thứ nhất, nó đổi vị trí của
chương~\ref{ch:chi-so-khong-do-do-trung-thuc} trong sách: bài neo không phải một
bằng chứng mượn từ một họ lời giải thích khác, nó là bằng chứng duy nhất đúng
loại mà sách có, và nó nằm ở phía không phải phía mà sáu chương quanh nó bàn tới.
Thứ hai, nó đổi cách đọc bốn bài của chương này: khảo sát 2026 gọi tên ba họ
dụng cụ mà không nói họ nào đã được kiểm, trong khi cả ba đều đã bị đem ra kiểm
và không họ nào đứng vững, nên chỗ hổng ở đây không phải chỗ hổng về bằng chứng
mà là chỗ hổng về việc bằng chứng đã có không tới được nơi cần tới. Thứ
ba, nó đổi câu hỏi mà chương~\ref{ch:khoang-trong-nghien-cuu} phải hỏi: không
phải \enquote{khi nào ai đó sẽ đo}, mà \enquote{vì sao một phép đo đã chạy rồi
lại không đổi được cách hai khảo sát của chương này viết về đối tượng của
chúng}.

\index{chuỗi suy luận!chỗ đặt chương trong sách}%
Chỗ mà sách đặt chương này cũng phải sửa theo, và đó là việc của chính mục này
chứ không của người đọc. Luận đề ở mục~\ref{sec:ch01-luan-de-cua-sach} gọi đường
thứ tư là chuyển câu hỏi sang chuỗi suy luận, và
mục~\ref{sec:ch16-ba-cai-gia} giao chương này đi bằng chữ ấy. Chữ \enquote{chuyển}
đúng ở chỗ nó nói hai phía là hai phía, và sai ở chỗ nó gán cho phía này vai
người nhận. Cái mà chương này nhận từ Phần~IV là một câu hỏi chưa có lời đáp;
cái nó gửi đi thì đã gửi từ chương~\ref{ch:chi-so-khong-do-do-trung-thuc}, tám
chương trước, và Phần~IV vẫn đang cầm.

\index{khoảng trống nghiên cứu!cái chương này để lại}%
Ba thứ đi tiếp từ chương này, và cả ba đi tới
chương~\ref{ch:khoang-trong-nghien-cuu}.

Thứ đi tiếp thứ nhất là chính câu về chiều, cùng với ba hệ quả của nó. Một
chương lấy giao của các phát biểu hạn chế cần biết rằng một trong các câu hỏi nó
sắp liệt kê đã có lời đáp ở một phía, và lời đáp ấy không lan sang được.

Thứ đi tiếp thứ hai là đề nghị bỏ độ trung thực để thay bằng
\tn{khả năng giám sát}{monitorability}, mà mục~\ref{sec:ch17-bai-neo-tro-ve}
thuật lại qua bài neo. Sách chưa đọc hai bài
đứng sau đề nghị ấy, nên nó đi sang chương sau ở dạng một dữ kiện cần xét chứ
không phải một lập trường đã cân.

Thứ đi tiếp thứ ba là điều kiện thứ năm đặt lên một dụng cụ đo tương lai, sau bốn
điều kiện mà mục~\ref{sec:ch13-tran-noi-gi}, mục~\ref{sec:ch14-dieu-chua-giai-quyet},
mục~\ref{sec:ch15-doc-duoc-dieu-khien-duoc} và mục~\ref{sec:ch16-ba-cai-gia} đã
đặt. Điều kiện ấy đến từ chỗ ba bài của chương này in một cái tên cho hai quan hệ
khác nhau: \textbf{một dụng cụ đo phải nói nó hỏi đường mà mô hình đã đi, hay hỏi
ảnh hưởng nhân quả của lời giải thích lên đáp án, vì đó là hai câu hỏi và
bảng~\ref{tab:ch01-bon-quan-he} xếp chúng vào hai dòng.} Điều kiện này khác bốn
điều kiện trước ở chỗ nó không đòi công bố thêm thứ gì; nó chỉ đòi gọi đúng tên
thứ đang đo, và chương này gặp ba bài không làm được việc ấy.

\index{chuỗi suy luận!bốn chỗ đứng của bốn bài|)}%
Chương~\ref{ch:khoang-trong-nghien-cuu} khép lại sách bằng cách lấy giao của các
phát biểu hạn chế mà mười sáu chương đã thu, và nó nhận từ chương này một thứ mà
mười sáu chương kia không đưa: một chỗ mà phép đo đã được làm.
